\chapter{OAK-D-LITE DEPTH CAMERA}
\label{ch:camera}

The OAK-D-Lite is the robot's only camera and its only sensor that can tell
a dog from a chair. It is a \emph{stereo} depth camera, not a time-of-flight
(ToF) camera: it measures depth by matching features between two image
sensors a known distance apart, rather than by timing emitted light. The
practical differences are that it needs texture in the scene (a blank wall
gives no depth), it works outdoors in sunlight where many ToF sensors
saturate, and its depth error grows with the square of range
\eqref{eq:depthres} rather than staying roughly constant. This chapter
describes the hardware, the geometry that turns pixels into metres, the
on-camera neural network, and the calibration state.

\section{HARDWARE}
\label{sec:oakhw}

\begin{table}[htbp]
\centering
\caption{OAK-D-Lite specifications (Luxonis datasheet)}
\label{tab:oak}
\small
\begin{tabularx}{\linewidth}{lX}
\toprule
Item & Specification \\
\midrule
Processor & Intel Movidius Myriad X VPU: 16 SHAVE vector cores, neural compute engine, hardware stereo-depth and warp engines; about 4 TOPS total, 1.4 TOPS for neural networks \\
Colour camera & Sony IMX214, 13 MP (4208$\times$3120), rolling shutter, autofocus (\SI{8}{cm} to infinity), DFOV/HFOV/VFOV \SI{81}{\degree}/\SI{69}{\degree}/\SI{54}{\degree} \\
Stereo pair & 2$\times$ OmniVision OV7251, 640$\times$480, global shutter, fixed focus, DFOV/HFOV/VFOV \SI{86}{\degree}/\SI{73}{\degree}/\SI{58}{\degree} \\
Stereo baseline & \SI{75}{mm} nominal; \SI{74.3}{mm} from this unit's calibration \\
Depth range & about \SI{0.2}{m} (extended disparity) or \SI{0.35}{m} (standard) to over \SI{10}{m}; useful to about \SI{8}{m} \\
Interface & USB-C, USB~3.2 Gen~1 capable; on the BeagleBone it runs at USB~2 \\
Power & \SIrange{4}{6}{W} typical at \SI{5}{V}, up to about \SI{1.5}{A} peaks; powered from the USB port only \\
Size and mass & \SI{91}{mm}$\times$\SI{28}{mm}$\times$\SI{17.5}{mm}, \SI{61}{g} \\
\bottomrule
\end{tabularx}
\end{table}

Table~\ref{tab:oak} gives the specifications. All image processing, stereo
matching and neural-network inference run on the Myriad~X; the BeagleBone
receives only detection results (label, confidence, bounding box and 3-D
position), so the camera adds almost no CPU load to the board. Its cost is
power: at \SIrange{0.8}{1.5}{A} it is the largest single load on the cape's
\SI{5}{V} rail (\S\ref{sec:powerbudget}).

\paragraph{Software.} The board runs depthai~2.19.1 for Python~3.7 (the
newest wheel that installs on the board's Debian). At USB~2 the link carries the detections and a preview
stream easily; full-resolution video would not fit.

\section{PINHOLE CAMERA MODEL}
\label{sec:pinhole}

Each sensor is modelled as a pinhole camera with lens distortion removed by
the calibration. A point $\bm P=(X,Y,Z)$ in the camera frame ($X$ right, $Y$
down, $Z$ forward along the optical axis) projects to pixel $(u,v)$ by
\begin{equation}\label{eq:pinhole}
\begin{bmatrix}u\\v\\1\end{bmatrix}
= \frac{1}{Z}\,
\underbrace{\begin{bmatrix}f_x & 0 & c_x\\ 0 & f_y & c_y\\ 0&0&1\end{bmatrix}}_{K}
\begin{bmatrix}X\\Y\\Z\end{bmatrix},
\end{equation}
with focal lengths $f_x,f_y$ in pixels and principal point $(c_x,c_y)$. For
an image $W$ pixels wide with horizontal field of view $\theta_h$,
\begin{equation}\label{eq:focal}
f_x=\frac{W/2}{\tan(\theta_h/2)} .
\end{equation}
For the mono cameras at $W=640$ and the datasheet $\theta_h=\SI{73}{\degree}$
this is $f\approx\SI{432}{px}$; the calibration measured
$f\approx\SI{450}{px}$ (Table~\ref{tab:oakcal}), and the calibrated value is
used from here on. Inverting \eqref{eq:pinhole} at a known depth $Z$
gives the lateral position
\begin{equation}\label{eq:backproject}
X=\frac{(u-c_x)\,Z}{f_x},\qquad Y=\frac{(v-c_y)\,Z}{f_y}.
\end{equation}

\section{STEREO DEPTH}
\label{sec:stereo}

\paragraph{Rectification.} Using the calibrated intrinsics and the rotation
and translation between the two mono cameras, the Myriad's warp engine
remaps both images so that their image planes are coplanar and their rows
are aligned. After rectification a point appears on the \emph{same row} in
both images, and the search for its match is one-dimensional along that
row. The calibration's epipolar error (how far matching points stray from a
common row) measures how well this worked; this unit's is 0.36~px.

\paragraph{Disparity and depth.} If a point appears at column $u_\ell$ in
the left image and $u_r$ in the right, its disparity is $d=u_\ell-u_r$, and
similar triangles give
\begin{equation}\label{eq:stereo}
Z=\frac{fB}{d},
\end{equation}
with baseline $B=\SI{0.0743}{m}$. On this camera $fB\approx
450\times0.0743=\SI{33.4}{px.m}$. The matcher is a census-transform cost
with semi-global aggregation \cite{hirschmuller2008stereo}, run in hardware
over a search range of 96 disparities (0--95), or 191 with extended
disparity. The largest disparity sets the minimum depth,
\begin{equation}
Z_{\min}=\frac{fB}{d_{\max}}=\frac{33.4}{95}\approx\SI{0.35}{m}
\quad\text{(standard)},\qquad
\frac{33.4}{190}\approx\SI{0.18}{m}\quad\text{(extended)}.
\end{equation}
The chase pipeline uses the standard range, so a dog closer than about
\SI{0.35}{m} has no valid depth and falls back to bearing-only
(\S\ref{sec:chasemeas}). The \SI{0.5}{m} stop distance keeps the robot outside
that zone.

\paragraph{Depth resolution.} Differentiating \eqref{eq:stereo}, a disparity
error $\delta d$ gives a depth error
\begin{equation}\label{eq:depthres}
\delta Z = \frac{Z^2}{fB}\,\delta d .
\end{equation}
Table~\ref{tab:depthres} evaluates it for one whole pixel of disparity. With
sub-pixel interpolation (3 fractional bits, so steps of $1/8$~px) the
quantization is eight times finer, though real matching noise is typically
a few tenths of a pixel.

\begin{table}[htbp]
\centering
\caption{Stereo depth resolution, $fB=\SI{33.4}{px.m}$, $\delta d=1$~px}
\label{tab:depthres}
\small
\begin{tabular}{rrr}
\toprule
Range $Z$ & Disparity $d$ (px) & $\delta Z$ \\
\midrule
\SI{0.5}{m} & 66.8 & \SI{0.007}{m} \\
\SI{1}{m} & 33.4 & \SI{0.030}{m} \\
\SI{2}{m} & 16.7 & \SI{0.12}{m} \\
\SI{3}{m} & 11.1 & \SI{0.27}{m} \\
\SI{5}{m} & 6.7 & \SI{0.75}{m} \\
\SI{8}{m} & 4.2 & \SI{1.9}{m} \\
\bottomrule
\end{tabular}
\end{table}

At the \SI{0.9}{m} follow distance the range error is a few centimetres; at
\SI{5}{m} it is most of a metre. The chase controller therefore needs its
range only to decide how fast to close in, and steers on bearing, which is
accurate at any range.

\paragraph{Depth aligned to colour.} The chase pipeline sets the stereo
node's depth alignment to the RGB camera, so each depth pixel is re-projected
into the colour camera's view: back-projected with \eqref{eq:backproject}
using the right mono camera's intrinsics, moved by the calibrated
mono-to-RGB extrinsics, and projected with the RGB intrinsics. The aligned
depth map is $640\times400$. This lets a bounding box found in the colour
image index the depth map directly.

\section{SPATIAL DETECTION}
\label{sec:spatial}

The \texttt{MobileNetSpatialDetectionNetwork} node combines a detector
running on the colour preview (Chapter~\ref{ch:chase} describes the
network) with the aligned depth map. For each detection with normalized box
$[x_{\min},x_{\max}]\times[y_{\min},y_{\max}]$:
\begin{enumerate}
\item The box is shrunk about its centre by the \emph{bounding-box scale
  factor} $s=0.5$ to a region of interest
  \begin{equation}
  \mathrm{ROI}=\Bigl[\bar x-\tfrac{s}{2}w,\ \bar x+\tfrac{s}{2}w\Bigr]\times
               \Bigl[\bar y-\tfrac{s}{2}h,\ \bar y+\tfrac{s}{2}h\Bigr],
  \end{equation}
  where $(\bar x,\bar y)$ is the box centre and $w,h$ its size. Shrinking
  keeps the region on the animal's body and off the background around its
  outline.
\item The depth pixels in the ROI that lie inside the thresholds
  $[\SI{100}{mm},\SI{8000}{mm}]$ are averaged to give $Z$. Pixels with no
  valid disparity (value 0) are excluded.
\item The ROI centre $(u,v)$ and $Z$ are back-projected with
  \eqref{eq:backproject} using the intrinsics of the aligned depth map,
  giving $(X,Y,Z)$ in millimetres. (DepthAI reports $Y$ positive
  \emph{up}, the opposite of \eqref{eq:pinhole}; the chase node uses only
  $X$ and $Z$.)
\end{enumerate}
If no pixel in the ROI is valid, $Z=0$ is reported and the chase node falls
back to a bearing from the box position alone.

\section{IMAGE GEOMETRY OF THE DETECTOR INPUT}
\label{sec:previewcrop}

The colour camera runs at 1080p ($1920\times1080$) and the network takes a
$300\times300$ \emph{preview}. With DepthAI's default
\texttt{setPreviewKeepAspectRatio(True)}, the square preview is made by
cropping the central $1080\times1080$ of the frame and scaling it, rather
than squashing the full width. The preview then spans a horizontal field of
view of
\begin{equation}\label{eq:cropfov}
\theta_{\mathrm{prev}}=2\arctan\!\Bigl(\frac{1080}{1920}\tan\frac{\SI{69}{\degree}}{2}\Bigr)
\approx\SI{42}{\degree},
\end{equation}
not the full \SI{69}{\degree}. (With the calibrated RGB focal length,
$f=\SI{1516}{px}$ at 1080p, the full frame spans \SI{64.7}{\degree} and the
crop \SI{39}{\degree}; the datasheet and calibration differ by a few degrees.) Two consequences:
\begin{itemize}
\item A dog at the edge of the colour frame, outside the central
  \SI{42}{\degree}, is not seen by the network at all. The detector's field
  of view is narrower than the camera's.
\item The chase node's original fallback bearing used
  \texttt{rgb\_hfov\_deg} $=68.8$ for the preview, which made it about 75\%
  too large. The depth-based bearing $\atantwo(X,Z)$ is unaffected because
  DepthAI maps the ROI back to the aligned depth frame with the correct
  geometry.
\end{itemize}
PR~\#3 sets \texttt{setPreviewKeepAspectRatio} explicitly from a
\texttt{preview\_keep\_aspect} parameter (default true: centre crop) and
computes the preview's field of view at start-up from the RGB intrinsics in
the camera's calibration, $\theta_p=2\arctan\bigl(w/(2f_x)\bigr)$ with
$w=1080$ for the crop or 1920 for the full width: about \SI{39}{\degree} with
this unit's calibration. Setting \texttt{preview\_keep\_aspect} false gives the
network the full width at the cost of a horizontally squashed image.
\influx{Merged in PR~\#3 but not yet run on the robot.}

\section{CALIBRATION}
\label{sec:oakcal}

The camera's EEPROM held no calibration when it arrived on the robot, and
DepthAI refuses to compute depth without one. It was recalibrated on the
Windows PC with the depthai~v2 calibration tool and a printed ChArUco board,
then flashed (2026-10-01). Results:
\begin{itemize}
\item Left--right epipolar error: 0.36~px, inside the \num{0.5} target in
  \file{docs/CALIBRATION.md}.
\item Baseline: \SI{7.43}{cm} against \SI{7.5}{cm} nominal.
\item RGB lens position (autofocus step at calibration): 72. The RGB camera
  must be at the same lens position for the RGB--depth alignment to be
  exact; autofocus moving the lens shifts the RGB intrinsics slightly.
\end{itemize}
Table~\ref{tab:oakcal} lists the calibrated intrinsics. The procedure is in
\file{docs/CALIBRATION.md}, and the calibration file itself is backed up in
\file{camera/pc-tools/}.

\begin{table}[htbp]
\centering
\caption{Calibrated intrinsics (from the flashed calibration file)}
\label{tab:oakcal}
\small
\begin{tabular}{lcccc}
\toprule
Camera & Calibrated at & $f_x$, $f_y$ (px) & $(c_x,c_y)$ (px) & Extrinsic translation \\
\midrule
RGB (IMX214) & $3840\times2160$ & 3032, 3031 & (1885, 1099) & reference \\
Left mono & $640\times480$ & 450, 450 & (335, 238) & $x=\SI{-7.43}{cm}$ to right mono \\
Right mono & $640\times480$ & 453, 452 & (338, 240) & $x=\SI{+3.74}{cm}$ to RGB \\
\bottomrule
\end{tabular}
\end{table}
\noindent Intrinsics scale with resolution: at 1080p the RGB focal length is
half the 4K value, \SI{1516}{px}. The PC must use depthai~v2:
the v3 tools write a calibration format the board's 2.19 library cannot read.

\section{MOUNTING AND FRAMES}
\label{sec:oakframes}

The camera faces forward on the front of the deck. Its optical frame
(\frm{oak\_rgb\_camera\_optical\_frame}, $Z$ forward) relates to the robot's
\frm{base\_link} (REP-103: $x$ forward, $y$ left, $z$ up) by
\begin{equation}\label{eq:cam2base}
\begin{bmatrix}x\\y\\z\end{bmatrix}_{\mathrm{base}}
=\begin{bmatrix}Z\\-X\\-Y\end{bmatrix}_{\mathrm{cam}}
+\bm t_{\mathrm{cam}},
\end{equation}
with $\bm t_{\mathrm{cam}}$ the camera's mounting offset. The chase node does
not use \eqref{eq:cam2base} today: it steers on the bearing in the camera
frame directly, which is equivalent while the camera sits on the robot's
centreline and the offset is small compared with the follow distance.
\influx{No static transform for the camera is published yet. The target
filter proposed in \S\ref{sec:targetkf} needs one.}
